\section{Year in Review: 2012}

\paragraph{Past}


WordPress has released a statistical summary for Gornahoor in 2012\footnote{\url{http://jetpack.me/annual-report/4840255/2012/}}, for those interested in taking a look. 170,000 page views are not very many in the grand scheme of things, yet for a blog that has no real identity, it is probably significant. The trend has been gradually increasing, but quality, not quantity, is the real goal. Note that three of the most frequent commenters are now missing, which is fine since they were the most disruptive and showed the least understanding. Recent comments, i.e., over the past several months, have been much more useful.

What is less interesting, and actually disappointing, is the quality of the search terms, which overwhelmingly come via Google. Despite having such extensive original material and new translations, Gornahoor shows up almost nowhere for the search on ``Evola''. However, we come up high for ``hermit crabs'' and periodically we get several hits for that; I always assume it is from a science assignment at some school somewhere. ``African ass'' comes up every now and then; I assume they are really disappointed. Keep this in mind next time you rely on Google to do your research for you.

\paragraph{Present}


The experience has been personally fruitful. Over the past year, I have heard from several readers via email or telephone about their own spiritual paths. Obviously, I cannot and do not tell them what to do, since Gornahoor is not evangelizing for any particular exoteric path. However, if our interaction has helped anyone achieve clarity, I am happy. Interestingly enough, everyone seems to choose a valid Western path with relatively recent ties to the Middle Ages; specifically, no Druids or Wiccans trust my judgment. I have had one disappointing failure, a revert back to neo-paganism. He pesters me with nonsense missives that prove nothing; he should use the comments or the forums to get better feedback.

\paragraph{Future}


I have written 638 posts and I am committed to 1001 before shutting Gornahoor down. That comes to several hundred thousand words, since I try to keep posts around 1000 words. At 3 or 4 posts per week, if I can maintain that, we will be done in around two years. One of the posts is usually a translation, although I am finding fewer interesting pieces to translate. Those by Charles Maurras and Guido di Giorgio are not very popular. Evola is getting stale. There is some old material, plus the large volumes of his philosophical system; these would not be worth the trouble to translate and it is difficult to pull passages out of context. Then, of course, there are the books on race. A year or so ago, I translated his book Synthesis of Racial Doctrine for private subscription. I could not get 50 people worldwide to sign up for a free translation, so I stopped. Perhaps I will start again, but charge for it this time; let me know.

There is a toll for this on my life and on my health, since each post takes effort; they are never mindless commentaries on the daily news. In addition, I am taking on commitments to write blogs for commercial enterprises, anonymously of course. These tasks are preventing me from some projects I have in mind which would require a sustained book length theme. Even worse, it is affecting my spiritual life and the time I have for my own contemplative reading. I have been using up capital accumulated earlier and that may be running dry.

Another contributor or two to share the load would be helpful. I can suggest topics. For example, someone may want to do a meditation on the labors of Hercules or the Romance of the Rose. We can be so limited to theory.

\paragraph{Initiation}


It is impossible to find a discussion of Tradition that is not obsessed with the notion of initiation. Is the issue of sacerdotal initiation vis-à-vis royal initiation truly a burning existential issue for anyone? If so, I would like to hear from the readers who are priests and kings to get their take on it.

I believe I have mentioned that I myself have been initiated into Mahayana Buddhism, since that option was available to me, plus the reasons I have mentioned several times already. I have recited the Heart Sutra, repeated the mantra endlessly, engaged in yantra meditation, and sat zazen. I know the sound of one hand clapping. Does that make me more reputable to anyone beyond the actual texts? Does anyone think to himself, `wow, that was written by a real initiate'? Yet it was, for what it's worth. If I bore you, then you could listen to the Beastie Boys instead; they are also initiates. Does anyone get the point or do I need to spell it all out?

\paragraph{Preparation}


That brings me to the real point, to wit, what is important is not initiation, but rather preparation for initiation. Otherwise, you will most likely end up with the spiritual equivalent of drinking wine out of a box. Obviously, being familiar with sound doctrine is necessary. Is your preparation including the following points, at a minimum?

\begin{itemize}


\item The World is Intelligible

\item The World is ordered

\item There are multiple states of being beyond the physical world

\item Man is a unity, yet composed of multiple sheaths or souls

\item And so on
\end{itemize}


Of course, this must be ``carnal knowledge'', that is, ideas you live out in the flesh. Your whole life centers around sound doctrine and you are not deluded by the snares of ordinary life. You remember you have chosen all that happens in your life or else it is the Will of God; better if you cannot tell them apart.

When you read religious and spiritual texts, you understand the symbolic meaning behind them (which does not necessarily deny their factual meaning.) You understand that esoteric teachings are not opposed to the exoteric and, most especially, they do not represent a competing set of beliefs. Rather, they indicate different states of consciousness (not necessarily states of being). Efforts must be made to recognize these states.

In other words, you cannot wait for an elusive initiation. It is a struggle, a search, and adventure; even if and when you find a good initiator, he will not show you everything. No, he will still make you struggle and search. There is no one right way for everyone; that is why the hero has 1000 faces.

\paragraph{Winding Down}


Although some few men (I hear from no women) have benefited, I am at this point not sure what else needs to be done. The original task was to serve as the catalyst for an elite who would go into the world. I have heard from students, and I pray they will someday be in an academic position where they can spread the ideals of Tradition. Others hope to influence youth education; still others plan to form a healthy family in today's unsound environment.

On the other hand, there is no shortage of ideas, if we know how to think, how to have patience, how to mind fast. The idea creates the thought which creates the word that leads to action. That is a law and it reveals the importance of good ideas.

Hence, in the time left, I want to make sure that Gornahoor actually exhibits, and not just documents, the principles. Probably no one has noticed that. For example, it should demonstrate organic rather than mechanical thinking. It should be about principles, not limited to opinions or facts. When the thought becomes the external word, it should adhere to the rules of grammar, the laws of logic, with the art of rhetoric.

Hence, the topics need to arise from intuitive awareness, not from technical planning. That is the way it will be, for those who continue to visit.

\paragraph{Year in Review Continued}

\begin{quotex}
The direction in which we may find awakening and liberation, the direction of liberation, is vertical and has nothing to do with the course of history.
\flright{\textsc{Julius Evola}, \emph{The Doctrine of Awakening}}
\end{quotex}


Rather than respond to comments, I'll address several issues here. As I've indicated several times, our goal at Gornahoor has been to prepare an intellectual elite. Over time, the audience has changed quite a bit to the point where there is now more of a common mind or at least a general agreement of principles. In post 1001, I will reveal everything.

The end of my involvement in Gornahoor is not necessarily the end of Gornahoor. As more writers become involved, it will develop a different identity. It could take some time, perhaps generations, for the elite to become effective; after all, the time of preparation began some 90 years ago. Some readers think I have the alchemical skill to translate spiritual ideas to the written text and they have benefited from it. That was the intent: to clear away false ideas to get to the core of Tradition. Once certain principles are understood, the ensuing choice and commitment becomes clear and necessary.

Ideally, there should be 100 new blogs bringing Tradition to the various walks of life, professions, arts, and crafts. So, I am not predicting an end, but rather a beginning.

\paragraph{Buddhism}


Readers must take care in what they read and not jump to conclusions that do not follow. I try to be precise. Furthermore, texts must be understood in the light of other texts. Reading in this manner is difficult; however, writing in such a way is even more difficult. I hope I have not failed completely. For example, I wrote:

\begin{quotex}
I myself have been initiated into Mahayana Buddhism, since that option was available to me.
\end{quotex}


Given what we have written about initiation and conversion, particularly in the case of Guenon, who can conclude from this that I am a Buddhist? It is obvious that Guenon and Evola have put great stock into the idea of ``initiation''; as the correspondence show, Evola pestered Guenon with that question even 25 years after their first exchange of letters. So I, too, took it quite seriously. I also took seriously Guenon's claim in the King of the World that the `lost word' would be found among the sages of Tibet.

Hence, I studied everything Tibetan, including occultists like Blavatsky, Nicholas Roerich, even Alice Bailey who claimed a Tibetan connection. I took the sangha vows and a few initiations from various Tibetan Lamas. Now that settled the problem of initiation for me so I could continue my studies of Tradition knowing I was an initiate. Was that necessary? Readers can decide for themselves. As for myself, I have been thinking lately ``no''. On the other hand, maybe personal instruction in forms of meditation and hearing doctrine from initiates really has made an effective difference in understanding; that can't be ruled out. In retrospect, it did lead to a change of direction in my life. However, I don't recommend this for all, especially those who have returned to the Western Tradition since they would have to confess that apostasy.

In any case, that did not make me a Buddhist; just an initiate in a valid Tradition. I left when I realized that my purposes had been served. In any case, Buddhism in America is in a sorry state. For example, I received a booklet about my dharma vows, in which there is a general prohibition against oral, anal, and masturbatory sex. Since a large number of American Buddhists are ex-Catholics who have told me they have moved beyond that, these prohibitions are mostly ignored. Even more curiously, a minor transgression against the vows can be repaired with a personal prayer, but a more serious one requires the retaking of the vows with a lama. This is identical to the Catholic distinction between venial and mortal sins.

What I liked about Tibet was that it was the last of the patriarchal theocratic societies, so it is still a good model. Unfortunately, no one else was interested in that sort of Tradition. Apparently, the Dalai Lama himself is thoroughly Westernized, despite his reactionary past if anything on the Internet about him can be believed. A blessing did come my way, from a Jewish convert to Buddhism who had a PhD in medieval history. She convinced me that the Medieval era was much more than the ``Dark Ages'', which gradually altered my view of that time. More detailed study led me to my current understanding.

As for the question about the Doctrine of Awakening, Buddhism has become a ``devotional faith'' by Aryan westerners. I oppose ``large'' judgments about movements that evolve over the centuries involving hundreds of millions of people. They give the impression of erudition while lacking all content. You can judge the Doctrine of Awakening by how well Evola extracts the various inner states of Buddhist practices. He does that well through his quotes and organization of key ideas. Nevertheless, that is merely academic exercise without making the efforts of actually achieving those states. They do not arise, I can assure you, from reading about them.

\paragraph{Priests and Kings Redux}


To return to this tedious theme once again, despite the many times it has come up, no one ever addresses anything specific in those posts. In the first edition of Revolt against the Modern World, Evola quoted a text from the Aitareya Brahmana and claimed that the King was the Sun and the Priest the Moon. Coomaraswamy, a Sanskrit scholar, pointed out Evola's error in a review. The text really said the exact opposite; the Priest recited the bridal vow that an Indian groom makes to his bride. That is the proper relationship between the Priest, as the Sun, and the King as the Moon. Evola pulled that reference from future editions of Revolt, but never changed his thesis. Draw your own conclusions about hiding contrary evidence. From now on, any comments on this issue will have to address that specific text from the Brahmana.

From a metaphysical point of view, it makes sense. The Priest is the link between the physical and the transcendent, i.e., the vertical. The King acts horizontally, although accepting the advice and the counsel of his priests. The vertical is Purusha, the unchanging, the masculine. Action is Prakriti; even in Tantra, Shakti is the feminine principle of action. To twist this to mean something else is a waste of intellectual capital. Together, the vertical and the horizontal form the sign of the cross, the sign of Triumph.

\paragraph{The Western Tradition}


So this brings us to the self-loathing that marks certain segments of the counter-Tradition in the West. They look everywhere for Tradition except for where it is. For us, that is the pre-schism Church, uniting the East and the West horizontally as well as the connection to the best of pagan Greece and Rome. We also need to mention the contribution of the Celts and the Franks to this synthesis.

This includes an esoteric Tradition that has been poorly preserved and also the Hermetic tradition that is also preserved and developed. Note that this is not the same as so-called ``Traditional'' elements in churches today, since they are usually suspicious of anything called esoteric, perennial, or Hermetic. Nevertheless, the esoteric teachings are not intended to replace the exoteric teachings; rather they are reserved for the few who interest themselves in such things.

In modern times, we have mentioned a particular trend that shows the initiation still exists in the West, as Guenon suspected. The Russian esoterist \textbf{G. O. Mebes} initiated both \textbf{Valentin Tomberg} and \textbf{Mouni Sadhu}\footnote{15 Jan 2021: Correction. Tomberg was not initiated by Mebes, but was in direct contact with some of his students.}. We have mentioned less about the latter. Sadhu was actually a Polish Catholic who lived in \textbf{Ramana Maharshi}'s ashram for a time before settling in Australia. Based on his own experiences and studies, he claims that the highest teachings of the West were not less than those of the East. He published several useful books of exercises on concentration, meditation, theurgy, and control of the mind. Some may find them helpful.

Otherwise, you are all on your own, some doing better than others.

\flrightit{Posted on 2013-01-13 by Cologero}

\begin{center}* * *\end{center}

\begin{footnotesize}
\begin{sffamily}
\texttt{Jason-Adam on 2013-01-14 at 02:56 said:}

The reason for the endless debate on priestly vs kingly intiation is the Evolian substitution of political activity for spiritual enlightenment. Too many people, especially from the Dugin camp today, think being a traditionalist means marching down the streets screaming death to the West and support for third world atheist socialist regimes rather than trying to spiritually leave the world and know oneself.

Your reference to a certain American rap musical group whose deceased lead singer was a Tibetan Buddhist initiate was made to illustrate I believe the need for study to make said initiation real, otherwise it remains virtual and of no benefit. Too many people today think initiation is like a magic formula that suddenly makes a person into an angel, they forget it requires hard work... ... ... .

I apologise if I go off course, but may I ask you, as a Buddhist, what your opinion is of Evola's Doctrine of Awakening ? I dont believe I've seen you comment much on it before and would like to know your thought on his hypothesis of Buddhism being an Aryan warrior anti-religious path that was later made into a devotional faith by the non-Aryan Asiatics... ... ..I find Evola's interpretation problematic but would like to know your thoughts.

\hfill

\texttt{Audrey on 2013-01-14 at 10:58 said:}

Shocking to hear no register for Evola, especially considering the quality and background width of material presented. But again, we live in a paradise of roses, or do we?

\hfill

\texttt{Audrey on 2013-01-14 at 11:00 said:}

As for shutting own, does that include deletion as well?

\hfill

\texttt{Audrey on 2013-01-14 at 11:19 said:}

Not sure if I count as a man, though.

\hfill

\texttt{Cassiodorus on 2013-01-14 at 14:17 said:}

My skepticism and scientism was effectively smashed to pieces by the writings of the Perennialist school. Taking their lead, Tradition brought me ``back'' to Catholic Christianity. Since then, I have invested myself in the works of the Church Father, but recently I have focused on the Angelic Doctor, Thomas Aquinas.

The more that I have learned about the Great Tradition , the more that I have come to doubt the transcedent unity of relgions. Is it not possible that the Christian revelation is something of a ``ground zero'' of the Divine descent and that the Trinity does, in fact, tell us something about the Divine on the level of the pure Absolute?

\hfill

\texttt{scardanelli on 2013-01-14 at 18:27 said:}

As I have mentioned, your writings have been extremely helpful to me as well. I was raised in Catholicism and have made the journey from progressivism, to radical environmentalism, to Heidegger and Nietzsche, to tradition. When I first stumbled upon this site I was but a ``traditionalist without a tradition,'' as most become after reading Evola. But little by little, I have come to see my native Catholicism in a newer (and older) light.

``And the end of all our exploring. Will be to arrive where we started. And know the place for the first time.''

This site has been a light in a darkening world, opening a way beyond the darkness. It would be selfish to ask you to continue your efforts further than you are able. As I seem to have come upon this site near the end, I hope that I can contribute in any small way that I am able.

\hfill

\texttt{Mihai on 2013-01-15 at 03:38 said:}

I don't know, Cassiodorus.

You really have to put yourself in the position of someone at the other end of the world. For them, Christianity is just some abstract religion practiced in a far away land. Just like for the common Christian man of the West, Taoism is some abstract asiatic philosophy and Hinduism is some weird practice of sitting cross-legged with a minimum of clothes on.

Of course, this is all after the advent of globalism, when contacts between different parts of the world became open to large masses of people, not only to some merchants and wanderers as was the case before. In those days, knowledge of far-away religions was non-existent.

CS Lewis wrote somewhere that ``the pole you're sitting next to seems a lot bigger than the other ones further down the road''. Of course, it wasn't in any context concerning the unity of religions, but it is a good metaphor for this case as well. The further geographically and mentally we are from a pole -- which in this case is a certain religions form- the smaller and more insignificant it seems to us. A reversal of perspectives also applies.

Though I do agree that for western man, Christianity is the most suitable.

\hfill

\texttt{Cassiodorus on 2013-01-15 at 11:33 said:}

Mihai,

I appreciate your thoughts on this matter. I totally resonate with your pole metaphor and it represents the view that I have always upheld. However, the problem that I have been wrestling with, of late, is the difficulty in reconciling Trinitarian theology with the central claim of the Perennial Philosophy. According to Frithjof Schuon, the Christian Trinity must be on the level of the ``relative Absolute''- the first determination of the Supreme Principle. But, to my lights, this runs contrary to the fundamental Christian claim that the Trinity tells us something of the Absolute at the highest degree of \emph{beyond} being. If you are interested, an excellent debate/exchange on this subject between the Christian platonist Robert Bolton and the Traditionalist Jose Segura can be found at sacredweb.com.

\hfill

\texttt{Avery on 2013-01-15 at 11:53 said:}

I have found that those who become influential within a particular tradition, like Seraphim Rose, credit traditionalism with awakening them to the reality of metaphysics, but only as a conduit to eventually finding higher truth as a member of an order.

\hfill

\texttt{Pickman on 2013-01-15 at 17:31 said:}

``The reason for the endless debate on priestly vs kingly intiation is the Evolian substitution of political activity for spiritual enlightenment.''

Excellent analysis. An angry nationalist will achieve more (against modernism) by physically pushing basic conservative principles while upholding the faith within a dogmatic manner. The internet has killed genuine action in many ways.

If one has learned over many years of struggle to develop a coherent character then action is inevitable and must be be taken as course of fighting dark forces in the name of Tradition. Since the noose narrows around the few who care it is of urgency that there are those to stand against the external crisis we find ourselves in.

``Too many people, especially from the Dugin camp today, think being a traditionalist means marching down the streets screaming death to the West and support for third world atheist socialist regimes rather than trying to spiritually leave the world and know oneself.''

Dugin is not a traditionalist and certainly given too much attention by European nationalists for his own worth. He is a Slavic national bolshevik at the root and never holds back his distaste for the ``West''. If his political theory (or scheming) was simply a cause of modern reaction then we could take a more favorable view but there are deeper layers of his native soul at work here. Basically his political theory is not for export.

``Your reference to a certain American rap musical group whose deceased lead singer was a Tibetan Buddhist initiate was made to illustrate I believe the need for study to make said initiation real, otherwise it remains virtual and of no benefit. Too many people today think initiation is like a magic formula that suddenly makes a person into an angel, they forget it requires hard work... ... ... .''

This is the new age spiritualist fallacy, there is a magic pill for anybody to become whatever they want from Brooklyn Jew to Buddhist (in the BB's sense)

at least nominally. Interior they are the at the will of the counter-initiation.

``I apologise if I go off course, but may I ask you, as a Buddhist, what your opinion is of Evola's Doctrine of Awakening ? I dont believe I've seen you comment much on it before and would like to know your thought on his hypothesis of Buddhism being an Aryan warrior anti-religious path that was later made into a devotional faith by the non-Aryan Asiatics... ... ..I find Evola's interpretation problematic but would like to know your thoughts.''

Evola's interpretation is more than reasonable as it goes to the original source. Some scholars have backed this analysis. Buddhism is suddenly appealing when stripped down from all the layers of sentimental muck endowed by profane explication. The eight-fold path is harsh and therefore will rarely be answered but by a few.

\hfill

\texttt{Pickman on 2013-01-15 at 17:47 said:}

You count as a conduit for male chromosomes and that is significant enough. You are the ``matter'', mother, the feminine lunar archetype that is the reflection of the mans genius and light. Men shape you, mold you, decide your fate and will, you do not and never owned yourself. This is the cosmic law weather you believe this rule or not and any deviance from it leads to chaos and hell. Just what women want unless they are tamed appropriately in their place.

Now listen up woman, you have no choice but to make way for the solar masculine light that is the cosmic law of the universe by divine right and holy imperial justice. It is what exists eternally within the metaphysical heavens as absolute and will by divine right return to manifest again on earth by divine rule of order. If the world crumbles, if priestesses rise from the bowl of modern puss to rule over the emasculated west, it will be a fleeting moment of extreme corruption against natural law, to which it in the end of time and beginning of a golden cycle it will be hammered down with divine justice. Gold outshines silver.

\hfill

\texttt{Saladin on 2013-01-15 at 23:01 said:}

Dear Cassiodorus,

I looked at the list of the contributors to sacredweb.com. Guenon was not listed and there was no mention Evola. This is enough (for me) to prove that this journal is run by ``Counter Initiation'' or at least has been co-opted by It. Bunch of safe academics, psudo-sufis, and ``scholarly oxen'' (with a few notable exceptions).

This may be heretical to some here but I think Frithjof Schuon and his followers have deviated from Guenon's teachings. This is a proof that ``Counter Tradition'' is very much attentive and extremely aggressive in frist identifying and then co-opting any resistance to Modernism!

\hfill

\texttt{Mihai on 2013-01-16 at 06:23 said:}

``but I think Frithjof Schuon and his followers have deviated from Guenon's teachings. ''

This is quite a strong statement. What makes you say that ? Do you say that he deviated completely, in the essentials, or in some particular respects ? And if Schuon ``deviated'' why hasn't Evola done that too, for example ?

Also, I find it a bit disturbing to say that Schuon deviated from Guenon, as if Guenon founded some new religion or cult, or something.

The only concern for me is who deviated and who conformed to the truth, not to another individual.

I am asking this (and I am waiting for your answer on this) because I am tired to see that even in this field, things have deviated to partisanship, skirmishes between followers of this or that writer and so on.

\hfill

\texttt{Scardanelli on 2013-01-16 at 10:52 said:}

Agreed. The search for purity of thought in others seems to be a stumbling block on the spiritual path. In the modern world, everything is mixed and confused and it must be enough to know and recognize counter traditional ideas as they arise without throwing the baby out with the bath water so to speak. We must learn to be discerning when it comes to ideas and not fall into the trap of the ``true believers'' of Guenon, Evola, Schuon, etc... Wasn't the division of those who seek to uphold tradition delineated as a tactic of the counter initiation in a recent post?

\hfill

\texttt{Cassiodorus on 2013-01-16 at 11:19 said:}

Schuon, a Guenonian deviation? I wouldn't sign my name to that. You can go to online articles at Sacred Web.

I'm struggling to create a link. Mihai, if you are so inclined, take a look at the essay, ``The Problematic of the Unity of Religions'' by the Perennialist Catholic, Jean Borella. I'd be curious to hear your reaction.

Best,

Cassiodorus

\hfill

\texttt{Logres on 2013-01-16 at 16:19 said:}

Cassiodorus, I haven't found the article yet, but here are some editorial remarks upon it:

\url{http://religioperennis.org/documents/Editorial/Issue5/editorialIII1.pdf}

Borella is certainly right to think that Christianity may be a ``revealer of revelation'', yet I don't know enough about Guenon to say whether or not he would have rejected this -- certainly Tomberg talks about the new religion as a hallower and purifier of what had already come.

\hfill

\texttt{Cassiodorus on 2013-01-16 at 20:27 said:}

Hi Logres,

You can just google that essay by Jean Borella- it can be found at sacredweb.com. Another essay of note that I think would be of great interest was posted recently by Charles Upton entitled, ``Were Rene Guenon and Frithjof Schuon Biased Against Love?'' in response to a charge made referring to Guenon as an ``eye without a body''.

\hfill

\texttt{Logres on 2013-01-16 at 21:44 said:}

This debate is immensely interesting, but we should be careful that we ``find what we seek'' : that is, if we approach the subject expecting to find deviations, then we shall find them. Jesus addresses this, when he tells the Pharisees that they have an evil eye. So, I think we should humbly approach Schuon and Guenon and even Borella's differences (they may simply be ill put) with an eye to discovering the unity. By their fruit, you will know them, and Guenon and Schuon need no defending on this score. I just started an article on the Marian roots of Schuon's spiritual lineage, which looks good:

\url{http://www.cutsinger.net/pdf/colorless\_light\_and\_pure\_air.pdf}

Cutsinger says right out that the article is ``not for anyone'', not for the skeptic.

\hfill

\texttt{Logres on 2013-01-16 at 21:48 said:}

Borella is a latecomer : that is, the Church has belatedly ``woke up'' to the perennial tradition, so there are problems with this situation that need to be kept in mind. As far as I know, even F. Rose was cautious in his critique of Guenon. He limited himself to saying that they were a waysign on the path of Truth, but not its fullness. But no one here, as far as I know, would want to commit to the same ecclesiastical views held (in their particular entirety) that F. Rose had, although I certainly wouldn't mind him praying for me. Cassiodorus, thanks, I did find it.

\hfill

\texttt{Mihai on 2013-01-17 at 04:47 said:}

This is something interesting. it seems that many people who started from Guenon and ended up embracing a Christian path have become skeptical of the unity of traditions. Borella and Fr. Seraphim Rose are perhaps the best known examples, but I see here other people too.

I don't know the cause of this. Is it indeed a matter of an inner superiority of the Christian revelation ? Or maybe the exclusivist spirit of the Church Fathers that is overwhelming in the long run ?

It would be interesting to hear Cutsinger's point on this, since he did not come to embrace this view.

\hfill

\texttt{Jason-Adam on 2013-01-17 at 08:28 said:}

Boris Mouravieff explained the matter quite clearly -- the other forms of the tradition were revelead to one group of humanity in particular whereas Christianity is universal and given to all men as a completion of the Tradition. Thus it is not contradictory for me to say all religions are derived from the one God but only the Catholic Church contains the fullness of truth.

\hfill

\texttt{Cologero on 2013-01-17 at 08:35 said:}

In a response to F. G. Galvao's claim that ``Schuon had realized Guenon's dreams'', Rene Guenon wrote back (2 December 1948):

\begin{quote}\footnotesize

I don't recall having used that expression but I have to recognize in it that it corresponds nicely to something true, for it is certain that without him, many things would have remained in the state of some sort of ``theoretical'' possibility and would not have been able to come to be realized se realiser. (My translation)
\end{quote}  

That is a rather strong endorsement, but perhaps something changed subsequently.

\hfill

\texttt{Mihai on 2013-01-18 at 06:43 said:}

``Thus it is not contradictory for me to say all religions are derived from the one God but only the Catholic Church contains the fullness of truth.''

For you it may not be contradictory, but for God it surely is.

What you just said implies that God, who is Truth, has been delivering to humanity something less than the truth. And anything less than the truth contains a degree of falsehood. Which is to say God is the author of falsehood around the globe.

Such explanations are not going to do.

\hfill

\texttt{Logres on 2013-01-18 at 06:55 said:}

The fullness of Truth can/could be ready for certain stages of humanity, if the makeup of the receiver determines (partially) the form and content of what is revealed; isn't this the whole point of the Old Testament? However, it is also true that ``not one jot or tittle shall pass away till all is fulfilled'' -- so that there is a paradox. A revelation can be (if the OT is any guide) fully complete, and totally the truth, containing within it the splendor of truth and fullness of beauty, and yet the fulfillment be such as to make it seem (if possible) lacking. I realize this analogy may not hold for other revelations, but the paradox exists. I think it has something to do with the ``double oath'', once by earth, once by heaven.

\hfill

\texttt{Cassiodorus on 2013-01-18 at 20:45 said:}

It seems to me that the question of Christianity and its ``compatibility'' with the Perennial Philosophy revolves around how we interpret the Incarnation, the Word Made \emph{Flesh. As I understand it, the Perennialist would draw a very important distinction between the statements ``Jesus is God '' and '' God is Jesus.'' Speaking for the Traditionalist School, the Perennialist Orthodox theologian, James Cutsinger, states the following:.}

\emph{``Christian Perennialists conclude that it is a mistake to confuse the uniqueness of the only-begotten and eternal Son of God with the alleged singularity of his historical manifestation in fist-century Palestine. Without denying that there is only one Son of God, or that he alone is the author of salvation, or that Jesus Christ is that Son, they contend that there are no Biblical or dogmatic grounds for supposing that this one Son has limited his saving work to his incarnate presence as Jesus. On the contrary, as St Athanasius and other fathers insisted, though the Word ``became flesh and dwelt among us'' (John 1:14), he was not confined by his body even during his earthly ministry.''}

Inclusivist Christians hold to the view that the the Third Person of the Trinity, the Holy Spirit, can be and is present in other world religions, rendering them salvific, but indirectly. The Perennialist position seems to be saying that it is Christ , the Logos, that is manifesting in different ways in other Traditions, in addition to Jesus of Nazareth . Can this view truly be reconciled with Christian orthodoxy?

\hfill

\texttt{Cologero on 2013-01-19 at 08:25 said:}

Cassiodorus, all things visible and invisible, were created through the Logos, obviously including different Traditions in some way. Guenon repeatedly claimed that all theological dogmas can be reframed in metaphysical terms; hence, discussions along these lines are well worth the effort. Now, Solovyov attempted to explain dogmas such as the Incarnation and the Trinity in terms of his metaphysical system, as well as those influenced by him. That should be worth a few studies at some time. Meanwhile, you could refer to Rama Coomaraswamy who tried to deal with this issue of exclusivity vis-a-vis perennialism. For example, what he writes about the \textit{Baptism of Desire}\footnote{\url{http://www.the-pope.com/bapodesc.html}} is of interest.

\hfill

\texttt{Graham on 2013-01-19 at 10:09 said:}

There is a difference between being ``derived from God'' and being ``delivered by God''.

I'm just pointing this out. I don't necessarily agree with JasonAdam's formulation.

But the idea that God would `hand-deliver' polytheistic religions is surely less tenable. See the first Commandment.

\hfill

\texttt{Cassiodorus on 2013-01-19 at 20:14 said:}

Your last comment in the forum was well said.

Vincit Omnia Veritas

\hfill

\texttt{Janus on 2013-01-16 at 01:39 said:}

Where in the west would you say one can find the pre-schism Church? In my journey, I have investigated both Catholicism and Orthodoxy. In Catholic churches today, one seems to be faced with a grim choice of a modernizing Church more concerned with condoms in Africa than transforming sacraments on the one hand, and a traditional Catholicism on the other which so often is filled with reactionary elements (in the sense of being hostile to anything perennial or esoteric, as you said). Orthodoxy seems to me to have far more promise, but as Fr. Rose found, hostility can come from this direction as well.

Personally, I stumbled upon Tradition as one not believing in God (a conclusion reached after the journey mentioned above). Investigating the Traditionalist writings, the idea of God as the ``face'' of the Divine provides some reconciliation here to me. As I once read it described, ``God is a mask. Who wears Him?'' The questions of Traditional metaphysics and personal practice are great ones indeed, even without considering the question of God in the theistic conception of Him. As such, Buddhism has increasingly seemed interesting to me, as it is not a theistic path. Would you say that I should look back to Catholicism or Orthodoxy despite this (as thoughtful advice, not guidance from a guru, of course)? One might have either ``a shaykh or Shaitan'' as ones' guide, but when shaykhs seem impossible to find in the remains of the the western Tradition, one really seems left to ones' own devices.

That said, I was reading earlier today Prince Charles' opening address to Sacred Web's 2006 conference. If the Prince of Wales himself is struggling with these questions, perhaps that provides some comfort that we are not so alone as it may appear, and hope that we may indeed see the ``end to all our exploring''. Thank you very much, Cologero, for the work you did here. I particularly enjoyed the Maurras translations, and would definitely like to see more if you are able. For Maurras and de Maistre alone, French is a language to learn. Best wishes, and I hope to see this blog and your work continue after you finish.

\hfill

\texttt{Avery on 2013-01-16 at 02:09 said:}

Janus,

Inherent in the juxtaposition of tradition and modernity is the recognition that things are at a very decayed state today. But I see the existence of this blog as a sign of hope. Guenon diagnosed the disease before the victim was entirely dead. After him there have not been many doctors with the intellectual ability to start treatment. This is the beginning of that era. It is natural, following the intellectual rigor of Guenon, to see decay all around you. Embrace it and correct it with compassion.

\hfill

\texttt{Brett Stevens on 2013-01-16 at 07:35 said:}

Tradition seems to me to be a truth derived from life itself that can be spoken in many languages or philosophical types. However, that doesn't mean it bends like a young tree in the wind. It is roughly the same thing because the world and cosmos are the same everywhere.

I hope you continue writing in whatever form suits you. From having written online for over 20 years, I recommend taking regular breaks anyway. You can simply deplete yourself of what you wish to communicate, and in doing so, neglect internal communication that helps you develop clarity of your own ideas.

\hfill

\texttt{Cologero on 2013-01-16 at 08:24 said:}

I'm afraid, Janus, that the only thing I can tell you is to not confuse a particular institution with the spiritual tradition itself, which is much larger. Also, practices which are suitable for monks are not always suitable for men in the world.

\hfill

\texttt{Visionsofglory14 on 2013-01-17 at 00:48 said:}

You are correct that actual initiation in the modern west might be helpful but is unnecessary and often detrimental if you get involved with the wrong institution or subgroup of an institution. In this Kali Yuga, all of the once carefully guarded sacred texts are easily available for everyone, so once a seeker has the general idea, initiation is perhaps even frivolous. That isn't to say guidance from a master won't be helpful. There's a ton to be misled by out there.

However, we can expect rot from top to bottom in these organizations, so self-education/direction is necessary. Former traditional bulwarks have completely lost their orientation even to the point of inversion of principle. Only in rare glimpses can the former spirit be witnessed. It only remains in disparate individuals. And you can often find that you have a better chance of communicating principle to any half-way decent person you meet than a member of an apparently traditional institution, so restoration is a red herring. Better to simply aim at self-development while communicating to those who cross your path than hoping for a restoration.

\hfill

\texttt{Visionsofglory14 on 2013-01-17 at 01:07 said:}

Doctrine of Awakening is an exemplary work. I don't know where I'd be without it. Evola and Nietzsche are two poles of the same spirit. Nietzsche alone forms a closed circuit as modern hands have neutered him, but Evola is like an outlet that makes spiritual combat possible again. Guenon is the opposite extreme from Nietzsche. Nietzsche=pure activity, Guenon=pure spirit. Evola is an attempt to mediate the two.

\hfill

\texttt{Cologero on 2013-01-17 at 07:39 said:}

How exactly do you ``mediate'' between ``pure spirit'' and ``pure activity''? You say Evola ``attempted'' to perform that mediation ... did he succeed or fail? If the former, how did he do it?


Guenon alone of the three understood the relationship between Purusha (spirit, masculine) and Prakriti (activity, feminine). Nietzsche seemed to intuit this when he wrote, ``when you go to a woman, bring a whip,'' but he probably never understood its true significance. To wit, activity must be guided by spirit, or, a fortiori, activity must be dominated by spirit.

For Nietzsche, there is no spirit and no pole, rather only power and appearance. That may explain why, judging from my facebook feed, Nietzsche appeals primarily to women and homosexuals.

\hfill

\texttt{Jason-Adam on 2013-01-17 at 08:24 said:}

1. I am still surprised that anyone still can admire a man like Nietzsche on the same level as a true man of wisdom such as Guenon or even Evola.

2. Cologero has done an amazing service in pointing us to Mebes, Tomberg and the others as they may be well the long sought after Christian path to transcendence that was occluded after the Middle Ages.

\hfill

\texttt{doctoroctagon on 2013-01-17 at 08:48 said:}

``our goal at Gornahoor has been to prepare an intellectual elite.''

I was put in mind of this: ``Never doubt that a small group of thoughtful, committed, citizens can change the world. Indeed, it is the only thing that ever has.''

Margaret Mead. Also my blog over at \url{http://www.golgonooza.blogspot.com} is going to have some more

Tradition and the arts oriented posts soon -- it might become one of those 100 blogs that you mention!

\hfill

\texttt{Visionsofglory14 on 2013-01-17 at 10:32 said:}

Evola, Guenon, and Nietzsche are united in that they recognize the decadence of modernity and that it must be destroyed so that life can regenerate. The greatest difference between Guenon and Evola is that Guenon sees the Kali Yuga as a blight, even though he recognizes as a necessary step in the return to the Golden Age. Nietzsche, on the other hand, teaches to enter the nihilistic phase with glee.

Nietzsche gave his bible Thus Spoke Zarathustra the subtitle ``A Book for All and None.'' A book that is meant for everyone will be vastly read, read at many different levels, yet clearly, as Nietzsche expected, to be greatly misunderstood. So Cologero `s attempt to defame Nietzsche as a tool for homosexuals and women is sidestepped by Nietzsche's quite obvious knowledge and intention that he would be used for destructive purposes by devious people.

Both Nietzsche and Guenon, however, do see envision a return to greater times. Unknowingly or not, they share this. Guenon's focus is on the preservation of tools that will bring back the golden age and be used by the future elite. Nietzsche, though, recognized that the negative polarity is necessary to divide mankind and make us hopelessly unequal, to the socialist's chagrin. Thus, with his books, he offers both high and low a new found freedom in power. For the high, Evola and Guenon may offer something: Ride the Tiger. For the low, this cannot mean anything more than the tools for their own destruction.

\hfill

\texttt{Cologero on 2013-01-17 at 12:44 said:}

Dear Delusions of Glory: First of all, I did not attempt to ``defame'' anything. I simply stated a fact. A fact is true or false; if you think it is false, then simply refute it. Others will make up their minds about who is ``deviously'' misusing Nietzsche.

Let me remind you, as well as everyone else, we are interested in discussing principles here, not in idle displays of pseudo-erudition of the type found in Spark Notes. So what exactly are the principles embraced by Guenon and Nietzsche? Obviously, Nietzsche claimed that the Will to Power is the fundamental principle of the world. That is, there is no transcendence, nothing hidden, just multiple manifestations of the will to power. There is nothing of Guenon in that, nothing of Tradition in that. Now to mediate --- something I asked you to explain, although you are not quite up to the task --- between Guenon and Nietzsche can mean nothing other than finding a higher principle that would incorporate and encompass both of them. So what is that mediating higher principle? Take your time answering ... You say Evola tried to do that. What he did was distinguish between the best and the worst of Nietzsche, as he called it, much like a boy how tries to pick the peas and carrots out from his beef stew. The question then becomes, what exactly is left of Nietzsche? Is Evola being devious and destructive by ``correcting'' Nietzsche? The same could be said of Coomaraswmay, which we documented here. Yes, we can rely on Nietzsche for his powerful poetry and heroic attitude, but is AKC also devious? Again, take your time, because we expect specific quotes and examples, not empty hand waving.

As for Nietzsche and glee ... does he truly demonstrate that in his life?

\hfill

\texttt{Mihai on 2013-01-18 at 06:47 said:}

If Nietzsche's end goals are taken seriously (instead of merely adopting some good points and means he presented) than he can only be a catalyst for one thing: the counter-tradition.

\hfill

\texttt{Visionsofglory14 on 2013-01-18 at 08:46 said:}

Explain Nietzsche's, Evola's, and Guenon's end goals, since you speak on such high authority.

\hfill

\texttt{Mihai on 2013-01-18 at 09:49 said:}

Explain the difference between Nietzsche's ``superman'' and Marx's ``end of history''- if there are any

(Hints: They are both relativistic, utopic, and they begin and end in the immanent, meaning that both are completely horizontal, lacking in any transcendence. )

\hfill

\texttt{Cologero on 2013-01-19 at 09:31 said:}

I have no obligation to explain anything to you, visions of glory. Rather, you have the obligation to properly prepare yourself before entering comments. Specifically, you need to refer to the actual content of posted texts and especially you need to respond to direct questions. The only authority comes from truth. If there is no truth, there is no authority. That is the fundamental issue and always the starting point. If you deny the existence of truth, or the ability to know it, then there is no independent objective authority. There is only the personal will to power to impose opinions through force.

\hfill

\texttt{andrewlmt on 2013-01-20 at 02:21 said:}

Reading the Church Fathers is good medicine. They are a way ``in''; I believe worthy guides for the Christian.

\hfill

\end{sffamily}
\end{footnotesize}



